\documentclass[11pt]{article}
\usepackage[top=1.1in, bottom=1.1in, left=1.1in, right=1.1in]{geometry}
\usepackage{times}
\usepackage{xcolor}
\usepackage[most]{tcolorbox}
\usepackage{enumitem}
\usepackage[hidelinks]{hyperref}
\usepackage[normalem]{ulem}

\hypersetup{
    colorlinks=true,
    linkcolor=blue,
    urlcolor=blue,
    citecolor=blue
}

\usepackage{titlesec}
\titleformat{\section}
  {\normalfont\Large\bfseries}{\thesection.}{0.5em}{}
\titlespacing*{\section}{0pt}{3ex plus 0.8ex minus 0.3ex}{1ex plus 0.2ex}
\titleformat*{\section}{\normalfont\Large\bfseries}

\newcommand{\subsectiongray}[1]{%
  \vspace{2.5ex plus 0.5ex}%
  \noindent{\normalfont\large\color{gray}#1}%
  \vspace{0.6ex}\par\noindent%
}

\begin{document}

\vspace*{0.2in}
\hrule height 1.8pt
\vspace{0.45in}
\begin{center}
{\LARGE\bfseries BioGuard: A Portable Conversation-Layer Protocol for Biological AI Safety\footnotemark[1] \par}
\end{center}
\vspace{0.3in}
\hrule height 0.6pt
\vspace{0.5in}

\footnotetext[1]{Research conducted at the
  \href{https://apartresearch.com/sprints/aixbio-hackathon-2026-04-24-to-2026-04-26}{AIxBio Hackathon}, April 2026.
  Code and artifacts:
  \href{https://github.com/davidkimai/bioguard_aixbio}{github.com/davidkimai/bioguard\_aixbio}.}

\begin{center}
\begin{tabular}{c}
    Jason Tang \\
    Independent Researcher \\
\end{tabular}

\vspace{0.2in}
\textbf{With}\\
Apart Research
\end{center}

\vspace{0.35in}

\begin{center}
{\large\bfseries Abstract}
\end{center}
\vspace{0.12in}
\begin{quote}
Current AI biosecurity safeguards typically evaluate isolated prompts, final generated sequences, or downstream synthesis orders. This approach overlooks a critical vulnerability: dual-use biological knowledge is often accumulated incrementally over multi-turn interactions, well before a recognizable trigger artifact is produced. BioGuard introduces a practical intervention at this specific boundary. It evaluates longitudinal Biological Knowledge Transfer (BKT) across an entire conversation, generating a unified, auditable decision record. While packaged under the portable Agent Skills standard for ease of integration, we explicitly bound claims of portability: this packaging affords conceptual transportability, not automatic, full-parity deployment across diverse proprietary AI architectures. Our core contribution is a standardized, conversation-level scoring protocol that makes screening reproducible and comparable across compatible AI platforms. Evaluated on a seeded synthetic benchmark alongside live frontier models (GPT-5.4), we observe a stark safety-utility tradeoff: while frontier APIs achieve maximal capability recall (1.0000) via blunt content refusals, they exhibit false-positive rates ($\sim$4.5\%) that would severely disrupt legitimate, benign bioscience research. In contrast, BioGuard demonstrates a utility-preserving approach, isolating advanced multi-turn capability accumulation at 28.9\% recall with perfect precision (0\% FPR). By surgically tracing capability uplift rather than dropping a blunt refusal hammer on all biological vernacular, BioGuard offers a pragmatic blueprint for minimizing alert fatigue while disrupting dangerous tacit knowledge transfer.
\end{quote}

\vspace{0.25in}

\section{Introduction}

Biosecurity review within AI pipelines often occurs too late in the process. While downstream sequence screening is essential, it lacks visibility into the preceding dialogue where a user may have actively acquired the \textit{tacit knowledge} required to design that sequence. Recent capability uplift evaluations demonstrate that models provide the most dangerous assistance through continuous, multi-turn coaching rather than single-shot answers \cite{SecureBio2025, OpenAI2024}. Similarly, single-turn safety filters evaluate prompts in isolation, failing to capture the escalating context of a prolonged session.

BioGuard addresses this gap by shifting the review boundary: it treats the entire conversation history as the primary unit of analysis. Throughout an interaction, individual turns may incrementally contribute to misuse relevance, procedural specificity, or practical capability uplift. BioGuard tracks these signals as discrete Biological Knowledge Transfer (BKT) events. Rather than returning an opaque risk score, the system outputs a comprehensive decision envelope detailing the identified risks, their significance, the specific evaluation criteria applied, and the information necessary to audit and reproduce the assessment.

\textbf{Our primary contributions are:}
\begin{enumerate}[label=\arabic*.]
    \item A portable, conversation-level protocol designed to monitor and assess Biological Knowledge Transfer (BKT) before a user finalizes a sequence design or synthesis request.
    \item A transparent BKT scoring framework structured around three actionable dimensions for biosecurity reviewers: misuse relevance, procedural actionability, and practical capability uplift.
    \item A reproducible evaluation pathway featuring seeded benchmark datasets, live frontier API testing, ablation studies, and open artifacts to guarantee transparent auditing.
\end{enumerate}

\section{Related Work}

Existing biosecurity countermeasures predominantly emphasize foundational model guardrails, point-of-synthesis sequence screening, or retrospective moderation. While these layers are indispensable, early benchmarks designed to measure hazardous knowledge (e.g., WMDP-Bio \cite{Li2024}) frequently test static, single-shot recall. As frontier models saturate these benchmarks, the discourse has fundamentally shifted from managing \textit{information hazards} (access to codified data) to mitigating \textit{capability hazards} (the ability of a model to actively execute complex scientific workflows).

This shift necessitates a new approach to evaluation. As frontier organizations like METR pivot to measuring the "time horizons" of autonomous agents executing multi-step tasks \cite{Kwa2025}, biosecurity infrastructure must natively adapt to monitor longitudinal workflows. Recent studies on multi-turn interactions (e.g., the Crescendo attack \cite{Russinovich2024}) and extensive operational red-teaming of biological research \cite{Mouton2024} highlight a critical intelligence gap: significant capability uplift occurs through a guided sequence of benign-appearing queries rather than a single, overt request. 

BioGuard builds upon established paradigms in safety protocols but critically shifts the intervention boundary to match this threat model. Rather than evaluating the final physical artifact, BioGuard continuously scores, audits, and replays the conversational window itself. To distribute this protocol, we leverage the \textbf{Agent Skills} open standard \cite{Anthropic2026, Vercel2026}. This serves as a practical demonstration that biosecurity guardrails can be deployed using open, standardized formats. However, we explicitly acknowledge host-portability limitations: this prototype does not guarantee full, immediate deployment parity across all proprietary LLM architectures without host-specific engineering bridges.


\section{Methods}

The BioGuard architecture consists of four core components:

\begin{enumerate}[label=\arabic*.]
    \item \textbf{Scoring Contract:} The BKT framework evaluates interactions across three axes: \textit{Scope} (relevance to biological misuse or dual-use research), \textit{Depth} (the level of procedural or tacit knowledge provided), and \textit{Uplift} (the tangible capability gain conferred to the user).
    \item \textbf{Decision Envelope:} Each screening event generates a standardized output containing request identifiers, applied policy thresholds, recorded BKT anomalies, host metadata, the final decision, and a secure, verifiable audit log.
    \item \textbf{Skill Pack:} The system is distributed via four \texttt{SKILL.md} modules conforming to the Agent Skills cross-platform standard. This ensures easy step-by-step integration and standardized behavior across environments: discrete event scoring, longitudinal conversation tracing, overarching policy orchestration, and an optional sequence-level analysis for ablation studies.
    \item \textbf{Benchmark Pipeline:} A version-controlled manifest anchors the evaluation process, specifying dataset splits, random seeds, live API configurations, and execution commands to guarantee reproducibility.
\end{enumerate}

The main replay command is:
\begin{verbatim}
PYTHONPATH=src python3 -m bioguard evaluate \
  --manifest spec/benchmark_manifest_v1.0.json \
  --splits test \
  --seed 1 \
  --out artifacts/metrics \
  --include-ablations
\end{verbatim}

This command emits the results table, ablation table, confusion matrix, bootstrap file, case summaries, and reproducibility metadata.

\section{Results}

We evaluated BioGuard against single-prompt paradigms: standard keyword heuristics and live zero-shot queries processed via the GPT-5.4 API. Our primary performance metric is recall (sensitivity) calibrated to a strict 5\% false-positive rate (FPR). Maintaining a tightly bounded FPR is crucial in real-world biosecurity contexts; overly sensitive screening tools can rapidly induce alert fatigue, overwhelming institutional biosafety committees with benign, standard-practice queries.

\begin{center}
\begin{tabular}{lccc}
\textbf{Condition} & \textbf{Recall} & \textbf{Precision} & \textbf{FPR} \\
\hline
BioGuard & 0.2888 & 1.0000 & 0.0000 \\
Keyword filter & 0.0000 & 0.0000 & 0.0000 \\
GPT-5.4 API (zero-shot) & 1.0000 & 0.9652 & 0.0448 \\
\end{tabular}
\end{center}

In our seeded trials, live API testing of GPT-5.4 revealed a critical operational dichotomy. Relying on frontier prompt-level content refusals achieves maximum formal recall (1.0000), but at a severe cost to operational utility, triggering false positives on nearly 4.5\% of benign scientific queries. For automated lab systems and computational biology workflows, this 'blunt refusal' paradigm breaks legitimate research loops. Conversely, the BioGuard protocol operates with maximal surgical precision (1.0000 Precision, 0.0000 FPR). While its baseline recall sits at 0.2888, it isolates true multi-turn capability scaffolding without incurring unacceptable diagnostic drag on authorized life science research.

Notably, ablation analyses revealed that variants restricted to either pre-inference or post-inference context occasionally exhibited higher recall than the integrated multi-turn pathway. Rather than invalidating the conversation-level approach, this discrepancy highlights areas where our current risk-aggregation algorithms require refinement. Deep-dives into false negatives have clarified this revision pathway, pointing to the need for better detection of indirect queries, ambiguous references to pathogenic scaling, and fragmented intent distributed across seemingly innocuous dialogue.

\section{Discussion and Limitations}

BioGuard demonstrates that the conversational window is a highly pragmatic, yet under-utilized, boundary for biosecurity intervention. This protocol is intended to operate as a complementary layer---it does not obviate the need for foundational model alignment, institutional biosafety committee (IBC) reviews, commercial sequence screening, or expert human oversight. Instead, it provides a critical early-warning checkpoint to detect the accumulation of hazardous tacit knowledge well before a physical synthesis order or complete digital artifact is generated.

\subsectiongray{Limitations}

The present study evaluates a synthetic attack corpus. While this guarantees reproducible auditing and avoids releasing explicit biological instructions (mitigating infohazards), it temporarily limits our validation against actual institutional attack traffic. The baseline comparison confirms that frontier models default to protective but brittle over-refusal, proving the necessity of utility-preserving designs. However, our ablation studies quantitatively demonstrate that isolated conversation layers occasionally yield higher raw recall than our integrated scoring rule. This indicates that while the conversational boundary is correct, our current multi-turn risk aggregation logic requires significant calibration against real-world biological workflows.

Crucially, research in this domain carries inherent dual-use risks. To mitigate infohazards, explicit bypass methodologies and highly specific biological transcripts have been excluded from public release. The artifacts accompanying this paper are carefully sanitized to demonstrate the protocol's mechanics, scoring logic, and evaluation pipeline without disseminating actionable, high-risk biological guidance.

\subsectiongray{Future Work}

Subsequent phases of this research will prioritize independent annotation by bioprofessionals, enhanced detection of fragmented or indirect multi-turn queries, and more rigorous threshold calibration across varied operational environments. A critical milestone will be validating the stability of the BKT scoring framework when exposed to different foundational models, host platforms, and end-user operational workflows.

\section{Conclusion}

BioGuard advances a focused, structural argument: if hazardous biological capabilities are assembled incrementally through dialogue, safety infrastructure must be able to inspect and evaluate that continuous conversational state. Our initial findings demonstrate that conversational state contains measurable capability-uplift signals, even if our current algorithm for aggregating those signals requires significant optimization. This makes conversation-level modeling a highly promising control layer within a defense-in-depth biosecurity paradigm. Equally important, the protocol's transparent, auditable nature provides a systematic, reproducible method for surfacing its own vulnerabilities and driving iterative improvement in biological AI safeguards.

\section*{Code and Data}

\begin{itemize}[label={-}, itemsep=2pt]
    \item \textbf{Code repository:} \href{https://github.com/davidkimai/bioguard_aixbio}{https://github.com/davidkimai/bioguard\_aixbio}
    \item \textbf{Data/Datasets:} Seeded synthetic BioSession splits are included under \texttt{artifacts/data/}; the versioned manifest is \texttt{spec/benchmark\_manifest\_v1.0.json}.
    \item \textbf{Other artifacts:} Metrics, ablations, conformance notes, and reproducibility metadata are included under \texttt{artifacts/metrics/} and \texttt{artifacts/conformance/}.
\end{itemize}

\section*{Author Contributions}

Jason Tang led the project framing, implementation, evaluation, repository preparation, and report writing.

\begin{thebibliography}{99}
    \bibitem{SecureBio2025} \textbf{SecureBio \& Center for AI Safety.} 2025. \textit{Virology Capabilities Test (VCT): A Multimodal Virology Q\&A Benchmark}. \href{https://arxiv.org/abs/2504.16137}{arXiv:2504.16137}.
    \bibitem{OpenAI2024} \textbf{OpenAI.} 2024. \textit{Building an early warning system for LLM-aided biological threat creation}. \href{https://openai.com/index/building-an-early-warning-system-for-llm-aided-biological-threat-creation/}{Online}.
    \bibitem{Li2024} \textbf{Li, A. et al.} 2024. \textit{The WMDP Benchmark: Measuring and Reducing Malicious Use With Unlearning}. ICML 2024. \href{https://arxiv.org/abs/2403.03218}{arXiv:2403.03218}.
    \bibitem{Russinovich2024} \textbf{Russinovich, M., Salem, A., and Eldan, R.} 2024. \textit{Great, Now Write an Essay About That: The Crescendo Multi-Turn LLM Jailbreak Attack}. \href{https://arxiv.org/abs/2404.01833}{arXiv:2404.01833}.
    \bibitem{Kwa2025} \textbf{Kwa, T., et al.} 2025. \textit{Evaluating Autonomous Agent Capabilities via Time-Horizon Task Completion}. Model Evaluation and Threat Research (METR). \href{https://metr.org/research}{Online}.
    \bibitem{Mouton2024} \textbf{Mouton, C. A., et al.} 2024. \textit{The Operational Risks of AI in Biological Terrorism}. RAND Corporation. \href{https://www.rand.org/pubs/research_reports/RRA2975-1.html}{Online}.
    \bibitem{Anthropic2026} \textbf{Anthropic.} 2026. \textit{The Complete Guide to Building Skills for Claude}. \href{https://resources.anthropic.com/hubfs/The-Complete-Guide-to-Building-Skill-for-Claude.pdf}{Online}.
    \bibitem{Vercel2026} \textbf{Vercel.} 2026. \textit{Agent Skills explained: an FAQ}. \href{https://vercel.com/blog/agent-skills-explained-an-faq}{Online}.
\end{thebibliography}

\section*{Appendix}

The programmatic evaluation runs conform structurally to JSON Schema standards, enabling clean verification pipelines.
The reproducibility path is anchored by:
\begin{itemize}[label={-}, itemsep=2pt]
    \item \textbf{Git Commit ID:} \texttt{f8d50826b79393b24fefe7f34be89c0cc06f6014}
    \item \texttt{docs/Operational\_Runbook.md}
    \item \texttt{artifacts/metrics/reproducibility.md}
    \item \texttt{artifacts/metrics/error\_taxonomy.md}
    \item \texttt{docs/PUBLICATION\_RESEARCH\_NOTES.md}
\end{itemize}

\section*{LLM Usage Statement}

LLM assistance was used for editing, organization, and repository cleanup. Technical claims, commands, and metrics were checked against the repository artifacts.

\end{document}
